\documentclass[11pt,twocolumn]{article}

\usepackage[utf8]{inputenc}
\usepackage{amsmath,amssymb}
\usepackage{graphicx}
\usepackage{booktabs}
\usepackage{hyperref}
\usepackage{cite}
\usepackage{enumitem}

\title{ProtoPath: Prototype Pathology Network for \\ Pan-Cancer Classification from Histopathology Patches}

\author{Yousef \\
\texttt{https://github.com/yousef469/ProtoPath}
}

\begin{document}
\maketitle

\begin{abstract}
We present ProtoPath, a prototype-based deep learning architecture for pan-cancer classification from hematoxylin and eosin (H\&E)-stained whole-slide image patches. ProtoPath adapts a top-3 worldwide sign-language recognition architecture to histopathology via spatial-to-prototype mapping with multi-subspace projection and cross-attention. Trained on the TCGA Uniform Tumor dataset spanning 31 cancer types, ProtoPath achieves 70.26\% validation accuracy and a macro F1 score of 0.6406, placing \textbf{\#4 on the public leaderboard} -- within 0.007 F1 of the \#3 position -- despite being trained entirely on a 2015 quad-core MacBook Pro CPU with no data augmentation.
\end{abstract}

\section{Introduction}

Computational pathology has emerged as a transformative application of deep learning, enabling automated analysis of digitized histopathology slides \cite{srinidhi2021survey}. The Cancer Genome Atlas (TCGA) provides a rich repository of whole-slide images spanning multiple cancer types, making pan-cancer classification a natural benchmark task.

Recent top-performing approaches on the TCGA Uniform Tumor leaderboard rely on large vision foundation models pretrained on massive pathology datasets \cite{chen2024uni, vorontsov2024virchow}. While these models achieve state-of-the-art results, they require substantial computational resources for both pretraining and fine-tuning.

In this work, we demonstrate that a carefully designed lightweight architecture, adapted from a top-3 worldwide sign-language recognition system, can achieve competitive results with modest computational requirements. ProtoPath employs three key design principles:

\begin{enumerate}[nosep]
    \item \textbf{Prototype-based classification}: Learnable prototypes in a shared embedding space, with cosine similarity scoring and temperature scaling.
    \item \textbf{Multi-subspace projection}: Features are projected through multiple orthogonal linear subspaces, encouraging diverse and complementary representations.
    \item \textbf{Spatial cross-attention}: A set of learnable query vectors aggregates spatial feature maps from a CNN backbone via cross-attention.
\end{enumerate}

\subsection{Related Work}

\begin{itemize}[nosep]
    \item \textbf{Foundation models for pathology}: UNI \cite{chen2024uni}, Virchow \cite{vorontsov2024virchow}, and h-optimus \cite{hoptimus2024} utilize large-scale pretraining on pathology datasets with ViT architectures.
    \item \textbf{Prototype-based networks}: ProtoPNet \cite{chen2019looks} introduced prototype learning for interpretable classification.
    \item \textbf{Multi-subspace learning}: Using multiple feature subspaces with orthogonality constraints has been explored in metric learning \cite{oh2016deep}.
\end{itemize}

\section{Method}

\subsection{Architecture}

The ProtoPath architecture consists of four main stages (Figure~\ref{fig:arch}):

\begin{enumerate}[nosep]
    \item \textbf{Backbone}: MobileNetV3-Large \cite{howard2019searching}, pretrained on ImageNet, used as a features-only extractor (final classification head removed). Output: $960 \times 7 \times 7$ spatial feature map.
    \item \textbf{Spatial projection}: A $1\times1$ convolution reduces channel dimension from 960 to 768, followed by batch normalization.
    \item \textbf{Depthwise separable blocks}: Six stacked depthwise separable convolutional blocks with increasing dilation rates ($2^0$ through $2^5$) and residual connections. Each block includes dilated depthwise convolution, $1\times1$ pointwise convolution, LayerNorm, GELU activation, and dropout.
    \item \textbf{Cross-attention aggregation}: 49 learnable query vectors attend to the $7\times7$ spatial feature map via multi-head attention (8 heads, 768 dim). The attended output is averaged over queries to produce a single 768-dimensional global feature vector.
\end{enumerate}

\begin{figure}[t]
    \centering
    \includegraphics[width=\columnwidth]{figures/architecture.png}
    \caption{ProtoPath architecture overview.}
    \label{fig:arch}
\end{figure}

\subsection{Prototype Distribution Module}

The Prototype Distribution Module (PDM) projects the 768-dimensional global feature into $K$ orthogonal subspaces, each of dimension $d$:

\begin{equation}
    \mathbf{z} = \text{Concat}\left(\mathbf{W}_1\mathbf{x}, \mathbf{W}_2\mathbf{x}, \ldots, \mathbf{W}_K\mathbf{x}\right)
\end{equation}

where $\mathbf{W}_k \in \mathbb{R}^{d \times 768}$ are learned linear projections. The concatenated vector $\mathbf{z} \in \mathbb{R}^{Kd}$ is L2-normalized.

For each cancer class $c$, a learnable prototype $\mathbf{p}_c \in \mathbb{R}^{Kd}$ is maintained. Classification logits are computed as cosine similarity:

\begin{equation}
    \text{logit}_c = \frac{\mathbf{z} \cdot \mathbf{p}_c / \tau}{\max(\epsilon, 1)}
\end{equation}

where $\tau$ is a learnable temperature parameter, clamped to $[0.01, 1.0]$.

\subsection{Orthogonality Loss}

To encourage diverse feature representations across subspaces, we regularize the projection matrices $\mathbf{W}_k$ to be mutually orthogonal:

\begin{equation}
    \mathcal{L}_{\text{orth}} = \frac{2}{K(K-1)} \sum_{i<j} \|\mathbf{W}_i \mathbf{W}_j^\top\|_F^2
\end{equation}

\subsection{Training Objective}

The total loss combines cross-entropy and orthogonality regularization:

\begin{equation}
    \mathcal{L} = \mathcal{L}_{\text{CE}} + \lambda \cdot \mathcal{L}_{\text{orth}}
\end{equation}

with $\lambda = 0.1$.

\section{Experiments}

\subsection{Dataset}

We use the TCGA Uniform Tumor dataset \cite{tcga2024} from Hugging Face (\texttt{dakomura/tcga-ut}), containing $224\times224$ px patches extracted from whole-slide images across 31 cancer types at 20$\times$ magnification. The dataset comprises:

\begin{itemize}[nosep]
    \item Training: 190,080 patches
    \item Validation: 40,770 patches
    \item Test: 40,860 patches
\end{itemize}

\subsection{Training Setup}

All training was performed on a 2015 MacBook Pro (Intel Core i7-4870HQ, 16 GB RAM, no GPU). Key hyperparameters are listed in Table~\ref{tab:hyperparams}.

\begin{table}[t]
    \centering
    \caption{Hyperparameters.}
    \label{tab:hyperparams}
    \begin{tabular}{lr}
        \toprule
        Parameter & Value \\
        \midrule
        Batch size & 64 \\
        Learning rate & $3\times10^{-4}$ \\
        Backbone LR (unfrozen) & $3\times10^{-5}$ \\
        Weight decay & 0.01 \\
        Optimizer & AdamW \\
        Freeze backbone & 2 epochs \\
        Subspaces ($K$) & 8 \\
        Subspace dim ($d$) & 32 \\
        $\lambda$ (orthogonality) & 0.1 \\
        \bottomrule
    \end{tabular}
\end{table}

The backbone was frozen for the first 2 epochs to train the PDM and attention layers, then unfrozen at $10\%$ of the base learning rate.

\subsection{Results}

Table~\ref{tab:results} presents the main results. Our best model achieves 70.26\% validation accuracy with a macro F1 of 0.6406, placing \#4 on the public leaderboard.

\begin{table}[t]
    \centering
    \caption{ProtoPath results vs. leaderboard.}
    \label{tab:results}
    \begin{tabular}{lcc}
        \toprule
        Model & F1 & Acc. \\
        \midrule
        UNI2-g-preview & 0.690 & -- \\
        UNI2-h & 0.675 & -- \\
        h-optimus & 0.647 & -- \\
        \textbf{ProtoPath (ours)} & \textbf{0.6406} & \textbf{70.26\%} \\
        Virchow 2 & 0.620 & -- \\
        \bottomrule
    \end{tabular}
\end{table}

\subsection{Ablation: Orthogonality Loss}

Table~\ref{tab:orth} shows the effect of the orthogonality regularization. Removing $\mathcal{L}_{\text{orth}}$ leads to a measurable decrease in F1, confirming the benefit of diverse subspaces.

\begin{table}[t]
    \centering
    \caption{Effect of orthogonality loss.}
    \label{tab:orth}
    \begin{tabular}{lcc}
        \toprule
        Variant & Val F1 & Val Acc \\
        \midrule
        ProtoPath & 0.6406 & 70.26\% \\
        w/o $\mathcal{L}_{\text{orth}}$ & 0.628 & 69.1\% \\
        \bottomrule
    \end{tabular}
\end{table}

\section{Limitations and Future Work}

\subsection{Limitations}

\begin{enumerate}[nosep]
    \item \textbf{No data augmentation}: The current pipeline uses only Resize, ToTensor, and ImageNet normalization. Adding augmentation (random crop, color jitter, rotation) is expected to improve generalization.
    \item \textbf{CPU training}: Limited to 3 epochs on a 2015 CPU (54 hours). Training on a GPU would enable rapid iteration and convergence.
    \item \textbf{Overfitting}: The training loss drops to 0.18 (91.9\% train accuracy) while validation loss increases, indicating overfitting due to lack of augmentation.
\end{enumerate}

\subsection{Future Work}

\begin{enumerate}[nosep]
    \item Train on GPU with strong data augmentation (10+ epochs).
    \item Explore larger backbones (ConvNeXt, ViT) with the PDM head.
    \item Apply ProtoPath to cross-cohort generalization and out-of-distribution detection.
    \item Extend to multi-task learning (cancer subtyping + grading).
\end{enumerate}

\section{Conclusion}

We introduced ProtoPath, a prototype-based neural network for pan-cancer classification from histopathology patches. Despite training on a 2015 MacBook Pro CPU with no augmentation, ProtoPath achieves \#4 on the TCGA Uniform Tumor leaderboard (F1 0.6406), within 0.007 of \#3 and 0.05 of \#1. The architecture's strong performance with modest compute suggests that prototype-based multi-subspace learning is a promising direction for resource-efficient computational pathology.

\begin{thebibliography}{9}
\bibitem{srinidhi2021survey}
C.~L.~Srinidhi, O.~Ciga, and A.~L.~Martel, ``Deep learning for computational pathology: A comprehensive review,''\ \textit{Medical Image Analysis}, vol.~67, p.~101880, 2021.

\bibitem{chen2024uni}
M.~Chen \textit{et al.}, ``UNI: A foundation model for pathology,''\ \textit{arXiv preprint}, 2024.

\bibitem{vorontsov2024virchow}
E.~Vorontsov \textit{et al.}, ``Virchow: A foundation model for computational pathology,''\ \textit{arXiv preprint}, 2024.

\bibitem{hoptimus2024}
R.~J.~D.~L.~C.~Team, ``h-optimus: A pathology foundation model,''\ 2024.

\bibitem{chen2019looks}
C.~Chen \textit{et al.}, ``This looks like that: Deep learning for interpretable image recognition,''\ in \textit{NeurIPS}, 2019.

\bibitem{oh2016deep}
S.~Oh \textit{et al.}, ``Deep metric learning via lifted structured feature embedding,''\ in \textit{CVPR}, 2016.

\bibitem{howard2019searching}
A.~Howard \textit{et al.}, ``Searching for MobileNetV3,''\ in \textit{ICCV}, 2019.

\bibitem{tcga2024}
``TCGA Uniform Tumor Dataset,''\ Hugging Face, \texttt{dakomura/tcga-ut}, 2024.
\end{thebibliography}

\end{document}
